\documentclass[12pt, fleqn]{article}
\usepackage[portuguese]{babel}
\usepackage{newpxtext, newpxmath}
\usepackage{amsmath}
\usepackage{titling}
\usepackage{enumitem}
\usepackage{geometry}
\usepackage{graphicx}
\usepackage{float}
\usepackage{listings}
\usepackage{xcolor}

\graphicspath{{./images/}}

\geometry{
    paper=a4paper,
    top=64pt,
    bottom=64pt,
    left=112pt,
    right=112pt
}

\lstset{
    basicstyle=\ttfamily\small,
    backgroundcolor=\color{gray!5},
    rulecolor=\color{gray!20},
    frame=single,
    framesep=16pt,
    xleftmargin=16pt,
    xrightmargin=16pt,
    captionpos=b
}

\renewcommand{\arraystretch}{1.2}

\setlength{\tabcolsep}{8pt}

\NewDocumentCommand{\course}{m}{\gdef\thecourse{#1}}
\NewDocumentCommand{\authors}{m}{\gdef\theauthors{#1}}
\NewDocumentCommand{\professor}{m}{\gdef\theprofessor{#1}}
\NewDocumentCommand{\semester}{m}{\gdef\thesemester{#1}}

\renewcommand{\maketitle}{
    \begin{titlepage}
        \begin{center}
            {\scshape Universidade Federal do Rio Grande do Sul\\Instituto de Informática} \\
            \vspace*{\fill}
            {\large \thecourse} \\[16pt]
            {\LARGE \thetitle} \\[36pt]
            \theauthors \\
            \vspace*{\fill}
            \theprofessor \\[2pt]
            \thesemester
        \end{center}
    \end{titlepage}
}

\course{Circuitos Digitais – INF01058}
\title{LAB 09 – Contador síncrono up/down}
\authors{
    Ana Cláudia Rodrigues — 343123 \\[2pt]
    Bruno Samuel Ardenghi Gonçalves — 550452
}
\professor{Prof. Mateus Grellert}
\semester{2025/1}

\begin{document}

\maketitle

\section{Introdução}

Neste laboratório desenvolvemos um registrador de 8 bits com função de carga paralela, armazenamento, incremento e decremento. Dado um sinal de 8 bits, um clock de 50MHz e um seletor de 2 bits, diminuimos o clock a 2.98Hz e, selecionado um modo, exibimos a saída do registrador em um visor de 2 dígitos hexadecimais de 7 segmentos. A divisão do clock e decodificação da saída são feitos utilizando circuitos implementados nos laboratórios anteriores. O projeto foi implementado e simulado pelo software Quartus, posteriormente sintetizado e testado em uma FPGA Cyclone III D0 (modelo EP3C16F484C6).

\subsection{Entradas}

\begin{itemize}
    \item \texttt{CLK}:\hspace{3pt} Clock com frequência $f$.
    \item \texttt{D[7..0]}:\hspace{3pt} Valor a ser carregado no registrador.
    \item \texttt{SEL[1..0]}:\hspace{3pt} Seletor de modo do regsistrador:
    \begin{itemize}[label=$\circ$]
        \item \texttt{00}:\hspace{3pt} Carga paralela ($Q_t \leftarrow D_t$)
        \item \texttt{01}:\hspace{3pt} Armazenamento ($Q_{t+1} \leftarrow Q_t$)
        \item \texttt{10}:\hspace{3pt} Incremento ($Q_{t+1} \leftarrow Q_t + 1$)
        \item \texttt{11}:\hspace{3pt} Decremento ($Q_{t+1} \leftarrow Q_t - 1$)
    \end{itemize}
\end{itemize}

\subsection{Saídas}

\begin{itemize}
    \item \texttt{HEX1[6..0]}:\hspace{3pt} Representação do bit mais significativo da saída para o visor de 7 segmentos.
    \item \texttt{HEX0[6..0]}:\hspace{3pt} Representação do bit menos significativo da saída para o visor de 7 segmentos.
\end{itemize}

\subsection{Módulos Principais}

\begin{itemize}
    \item \textbf{}{Contador síncrono de 8 bits}: O núcleo do projeto, que contém a lógica para as quatro operações.
    \item \textbf{Dividor de clock}: Módulo divisor de frequência para reduzir a velocidade do clock para visualização.
    \item \textbf{Decodificador BCD-7seg (x2)}: Módulos que convertem a saída binária de 4 bits para o código de 7 segmentos do display.
\end{itemize}


\section{Resultados}

\begin{table}[H]
    \centering
    \begin{tabular}{|l | c|}
        \hline
        \textbf{Total de elementos lógicos} & 24 \\
        \hline
        \textbf{Total de pinos} & 25 \\
        \hline
        \textbf{Área total} & 48 \\
        \hline
        \textbf{Total de registradores} & 8 \\
        \hline
    \end{tabular}
    \caption{Consumo de recursos da FPGA em área}
\end{table}

\begin{table}[H]
    \centering
    \begin{tabular}{|l | c|}
        \hline
        \textbf{Frequencia máxima} & 548.85 MHz \\
        \hline
        \textbf{Atraso crítico} & 1.753 \\
        \hline
    \end{tabular}
    \caption{Frequência máxima de operação e atraso crítico (\textit{slow}, 85°C)}
\end{table}

\section{Capturas de tela}

\begin{figure}[H]
    \includegraphics[width=\textwidth]{netlist.png}
    \caption{Captura de tela do netlist}
\end{figure}

\begin{figure}[H]
    \includegraphics[width=\textwidth]{waveform.png}
    \caption{Captura de tela da simulação em forma de onda com atraso}
\end{figure}

\section{Arquivo de restrições}

\begin{lstlisting}[caption=constraints.sdc]
create_clock -name clk -period 20 [get_ports CLK]

set_input_delay -clock clk -max 0.2 [all_inputs]
set_input_delay -clock clk -min 0.01 [all_inputs]
set_output_delay -clock clk -max 0.2 [all_outputs]
set_output_delay -clock clk -min 0.01 [all_outputs]  
\end{lstlisting}

\end{document}
